\documentclass{article}
% Source: Topic_01_Teacher_Edition.pdf, PDF pages 38-49 (printed pages 23A-30B)
\usepackage{amsmath}
\usepackage{amssymb}
\usepackage{graphicx}
\usepackage{tikz}
\usepackage{pgfplots}
\pgfplotsset{compat=1.18}
\usepackage{booktabs}
\usepackage{xcolor}

\newenvironment{lesson-meta}{}{}        % lesson code, title, objective, EQ, vocab
\newenvironment{item-analysis}{}{}      % Savvas's Example × DOK table
\newenvironment{model-discuss}{}{}      % Launch opener
\newenvironment{concept-box}{}{}        % in-lesson concept reference
\newenvironment{concept-summary}{}{}    % end-of-lesson summary
\newenvironment{example}[1]{}{}         % #1 = example number
\newenvironment{tryit}[1]{}{}           % #1 = tryit number
\newenvironment{practice}[2]{}{}        % #1 = practice number, #2 = DOK
\newenvironment{te-addendum}[2]{}{}     % #1 = anchor example, #2 = type
\newcommand{\answer}[1]{}               % answer key, inside any env
\newcommand{\placeholder}[2]{}          % #1 = type, #2 = description
\newcommand{\uncertain}[2]{#1}          % #1 = best guess, #2 = alternatives
\newcommand{\te}[1]{}                   % teacher-only inline annotation

\begin{document}


% Source page 38: 23A Lesson Overview
\begin{lesson-meta}
lesson: 1-3
title: Piecewise-Defined Functions
standards: none printed
practices: none printed
objective: I can graph and interpret piecewise-defined functions.

essential question: How do you model a situation in which different functions describe specific intervals of its domain?

\textbf{VOCABULARY}
\item piecewise-defined function
\item step function
\textbf{TOPIC VOCABULARY}
\te{No topic vocabulary list or numbered standards/practice codes are printed. I CAN statement and Essential Question are on PDF pages 39-40.}
\begin{tabular}{ll}
Lesson Code & 1-3 \\
Title & Piecewise-Defined Functions \\
Standards & none printed \\
I CAN Objective & I can graph and interpret piecewise-defined functions. \\
Essential Question & How do you model a situation in which different functions describe specific intervals of its domain? \\
Lesson Vocabulary & piecewise-defined function; step function \\
Topic Vocabulary & none printed \\
\end{tabular}
\end{lesson-meta}
\begin{te-addendum}{0}{GeneralizeETP}
\textbf{Lesson Overview: FOCUS}
\textbf{Objective} Students will be able to:
\item Create and graph piecewise-defined functions, including absolute value functions and step functions.
\item Create and use a piecewise-defined function from real-world data.
\item Write a piecewise-defined rule from a graph.
\textbf{Essential Understanding}
A piecewise-defined function is used to model situations in which there are different rules for different parts of the domain of the function. Absolute value functions and step functions are two types of piecewise-defined functions.
\textbf{COHERENCE}
Previously in earlier courses, students:
\item Created and graphed absolute value functions.
In this lesson, students:
\item Understand that an absolute value function is a piecewise-defined function.
\item Create and graph piecewise-defined functions including absolute value functions and step functions.
Later in this topic, students will:
\item Write the general rule for an arithmetic sequence as a piecewise-defined function.
\textbf{RIGOR}
This lesson emphasizes a blend of conceptual understanding and application.
\item Students understand that a piecewise-defined function has unique rules for defined parts of its domain.
\item Students apply their understanding of piecewise-defined functions to interpret the relationship between a function's variables within a given context, such as the relationship between pay rate and number of hours worked.
\end{te-addendum}
\begin{te-addendum}{0}{ELLAddendum}
\textbf{Vocabulary Builder}
Glossary is available at SavvasRealize.com.
\textbf{REVIEW VOCABULARY}
\begin{tabular}{ll}
English & Spanish \\
absolute value function & función de valor absolute \\
domain & domini \\
\end{tabular}
\textbf{NEW VOCABULARY}
\begin{tabular}{ll}
piecewise-defined function & función de fragmentos \\
step function & función escalón \\
\end{tabular}
\textbf{VOCABULARY ACTIVITY}
Have students recall their understanding of absolute value functions.
Q: Why does an absolute value function have more than one solution?
Discuss the fact that absolute value functions have two rules, one for when the value inside the absolute value symbols is negative, and one for when it is positive. Explain that there is a different rule for two different parts or pieces of the domain. Any function that has different rules for different parts or "pieces" of its domain is called a piecewise-defined function.
\end{te-addendum}
\begin{te-addendum}{0}{LearnTogether}
\textbf{Student Companion}
Students can do their work for the lesson in their Student Companion or on Savvas Realize.
% FIG: 38 963 787 1116 970
\placeholder{illustration}{Student Companion cover with colored loops on dark background, Student Companion and enVision Algebra 2 titles.}
\end{te-addendum}
\begin{te-addendum}{0}{GeneralizeETP}
\textbf{Mathematics Overview}
Students model situations with different rules for different parts of the domain. They represent and graph piecewise-defined functions, including absolute value functions and step functions.
\textbf{Applying Math Practices}
\textbf{Attend to Precision}
Students communicate precisely and use clear mathematical language when they explain the steps that need to be followed when graphing a piecewise-defined function.
\textbf{Look For and Make Use of Structure}
Students look for overall structure and pattern in mathematics when they identify that the linear pieces of a piecewise-defined function are increasing when the slope is positive and decreasing when the slope is negative.
\te{SavvasRealize.com: Program Resources.}
\end{te-addendum}

% Source page 39: 23B STEP 1 Explore
\begin{te-addendum}{0}{LearnTogether}
\textbf{MODEL \& DISCUSS}
\textbf{INSTRUCTIONAL FOCUS} Students explore two different types of graphs, one linear and one with gaps. Students understand that graphs can behave differently depending on the situation. This leads to students to make a connection between linear graphs and graphs of piecewise-defined functions.
\textbf{STUDENT COMPANION} Students can complete the Model \& Discuss activity in their Student Companion.
\end{te-addendum}
\begin{te-addendum}{0}{GeneralizeETP}
\textbf{Before—WHOLE CLASS}
\textbf{Implement Tasks that Promote Reasoning and Problem Solving ETP}
Q: How do stores A and B compare?
\answer{Store A sells guitar strings individually, and store B sells guitar strings in packages of 4.}
\textbf{During—SMALL GROUP}
\textbf{Support Productive Struggle in Learning Mathematics ETP}
Q: How do the stores' graphs represent the cost of the guitar strings?
\answer{Store A has a steady increase of 6 additional dollars for each guitar string. Store B sells 4 guitar strings for \$20, so the graph is constant at certain y-values, such as 20 when (x) is between 1 and 4, and then jumps up \$20 and is constant again.}
\textbf{For Early Finishers}
Q: How can you represent the stores' guitar string prices as equations?
\answer{Store A: (y=6x); Store B: (y=20) for (0<x\leq4), (y=40) for (4<x\leq8), (y=60) for (8<x\leq12), (y=80) for (12<x\leq16), (y=100) for (16<x\leq20).}
\textbf{After—WHOLE CLASS}
\textbf{Facilitate Meaningful Mathematical Discourse ETP}
Q: What is different about the shapes of the two graphs?
\answer{Graph A has a linear shape, while graph B has multiple horizontal or constant line segments (step-function).}
\end{te-addendum}
\begin{te-addendum}{0}{HabitsOfMind}
\textbf{Use with MODEL \& DISCUSS}
Q: Communicate Precisely Why do you use dots rather than line segments to graph these two functions?
\answer{You use dots to show whole numbers of strings only. It does not make sense to talk about buying a fractional part of a guitar string.}
\end{te-addendum}
\begin{model-discuss}
\textbf{MODEL \& DISCUSS}
A music teacher needs to buy guitar strings for her class. At Store A, she can buy a single pack of strings. At Store B, she can buy a bundle of 4 packs of strings.
% FIG: 39 958 638 1132 743
\placeholder{illustration}{Two packages labeled Guitar Strings: bundle of four at left with price tag 20 dollars, single package at right with price tag 6 dollars; loose coiled string below.}
A. Make graphs that show the income each store receives if the teacher needs 1--20 packs of guitar strings.
\answer{A. See graphs.}
% FIG: 39 679 946 909 1087
\placeholder{graph}{Sample Student Work Store A: black filled discrete points (n,6n) for integers n=1 through 20. x-axis Number of Strings, labels 0,2,4,...,20 and grid spacing 1; y-axis Income (dollars), labels 0,20,40,...,120 and grid spacing 10. Arrowed x/y axes; window [0,20] by [0,120]; no connecting lines.}
% FIG: 39 914 946 1142 1087
\placeholder{graph}{Sample Student Work Store B: black filled discrete points at x=1,2,3,4 with y=20; x=5,6,7,8 with y=40; x=9,10,11,12 with y=60; x=13,14,15,16 with y=80; x=17,18,19,20 with y=100. x-axis Number of Strings, labels 0,2,...,20, grid 1; y-axis Income (dollars), labels 0,20,...,120, grid 10. Arrowed x/y axes; no connecting lines.}
B. Describe the shape of the graph for store A. Describe the shape of the graph for store B. Why are the graphs different?
\answer{B. The graph for Store A is a straight line; the graph for Store B jumps between flat sections; they are different because Store B sells in packages, not individually like Store A.}
C. Communicate Precisely Compare the graphs for stores A and B. For what numbers of guitar strings is it cheaper to buy from store B? Explain how you know.
\answer{C. It is cheaper to buy from Store B if you are purchasing 4, 7, 8, 11, 12, or 14--20 guitar strings. You can compare the points on the graphs for the number of strings and find which one has a lower price.}
\te{The printed digital launch repeats these prompts, adding "Use the pen tool to make graphs" in A. It displays Go back, ESSENTIAL QUESTION, SOLUTION, Enter your answer., and 1 of 1. The page link says "Critique \& Explain is available at SavvasRealize.com." The printed instructional paragraph includes "leads to students to."}
% FIG: 39 847 253 992 469
\placeholder{illustration}{Digital launch repeats the two Guitar Strings packages and 20-dollar/6-dollar price tags above a blank square graph grid. Grid has no axis labels, tick numbers, or plotted marks; pen-tool palette at left and 1 of 1 control at bottom.}
\end{model-discuss}


% Source page 40: 23 STEP 2 Understand & Apply
\begin{te-addendum}{0}{LearnTogether}
\textbf{INTRODUCE THE ESSENTIAL QUESTION}
\textbf{Establish Mathematics Goals to Focus Learning ETP}
Introduce students to piecewise-defined functions. Explain that graphs can behave differently over different domains and may not be continuous.
\textbf{ESSENTIAL QUESTION} How do you model a situation in which different functions describe specific intervals of its domain?
\te{Examples are available at SavvasRealize.com.}
\end{te-addendum}
\begin{te-addendum}{1}{GeneralizeETP}
\textbf{EXAMPLE 1 Model With a Piecewise-Defined Function}
\textbf{Build Procedural Fluency From Conceptual Understanding ETP}
Q: Why can a function have more than one rule?
\answer{Pieces of a graph can have more than one rule for different parts of the domain. In this example, there are two different rules—one when Alani has worked 40 hours or less, and one for more than 40 hours.}
Q: What does the equation (P(x)=480+18(x-40)) represent in this situation?
\answer{In this equation, (x) represents the number of hours worked. Alani makes \$12/h for the first 40 hours, which equals \$480, and (18(x-40)) represents the additional hours Alani works after 40 hours at \$18/h. The total represents Alani's weekly earnings.}
\end{te-addendum}
\begin{example}{1}
\textbf{Model With a Piecewise-Defined Function}
\textbf{CONCEPTUAL UNDERSTANDING}
Alani has a summer job as a lifeguard. She makes \$12/hour for up to 40 hours each week. If she works more than 40 hours, she makes 1.5 times her hourly pay, or \$18/hour, for each hour over 40 hours. How could you make a graph and write a function that shows Alani's weekly earnings based on the number of hours she worked?
Step 1 Make a table of values and a graph.
\begin{tabular}{ll}
Hours Worked & Pay \\
20 & 240 \\
25 & 300 \\
30 & 360 \\
35 & 420 \\
40 & 480 \\
45 & 570 \\
50 & 660 \\
55 & 750 \\
\end{tabular}
% FIG: 40 933 338 1060 504
\placeholder{graph}{Alani's pay: red segment P=12x from (0,0) through filled points (20,240),(25,300),(30,360),(35,420),(40,480), followed by blue ray P=18x-240 through filled (45,570),(50,660),(55,750) with upper-right arrow. x-axis Hours Worked labeled 0,20,40,60, grid 10; y-axis Pay (dollars) labeled 0,200,400,600,800, grid 100. Arrowed axes; visible [0,60] by [0,900].}
\te{Callout: When (x>40), Alani's pay is (P(x)=(\$12)(40)+\$18(x-40)), or (P(x)=18x-240).}
\te{Callout: When (0\leq x\leq40), Alani's pay is (P(x)=12x).}
Step 2 Notice that the plot contains two linear segments with a slope that changes slightly at (x=40). A function that has different rules for different parts of its domain is called a piecewise-defined function.
Step 3 Write an equation for each piece of the graph.
(P(x)=\begin{cases}12x,&0\leq x\leq40\\18x-240,&x>40\end{cases})
\te{USE STRUCTURE: This notation is used for piecewise-defined functions to indicate the different functions at different intervals of the domain.}
\answer{(P(x)=\begin{cases}12x,&0\leq x\leq40\\18x-240,&x>40\end{cases})}
\end{example}
\begin{te-addendum}{1}{PurposefulQuestions}
\textbf{Additional Examples—ADDITIONAL EXAMPLE 1}
The table shows how much a car rental costs per day. Use the table to write and graph the function for the car rental cost based on the number of days rented. Explain.
\begin{tabular}{ll}
\# of Rental Days & Cost \\
1 & 40 \\
2 & 80 \\
3 & 120 \\
4 & 155 \\
5 & 190 \\
6 & 225 \\
7 & 260 \\
\end{tabular}
For the first 3 days, the car rental costs \$40 per day. After renting for 3 days, the rental cost is decreased to \$35 per day.
The equation for each piece of the graph is:
\answer{(f(x)=\begin{cases}40x,&\text{if }0\leq x\leq3\\35(x-3)+120,&\text{if }x>3\end{cases})}
% FIG: 40 167 1028 291 1147
\placeholder{graph}{Additional Example 1 discrete rental graph: orange filled points (1,40),(2,80),(3,120) with orange label f(x)=40x; blue filled (4,155),(5,190),(6,225),(7,260), blue label f(x)=35(x-3)+120. x-axis Number of Rental Days labels 0,2,4,6,8,10, unit grid; y-axis Cost labels 0,50,100,150,200,250,300, grid 25. Arrowed axes, no connecting lines.}
\textbf{Example 1} Help students use information given in a table to write and graph a piecewise-defined function.
Q: Is there more than one way to represent the function?
\answer{Yes, when (x>3), the function can also be represented as ((35x-105)+120).}
\end{te-addendum}
\begin{te-addendum}{2}{PurposefulQuestions}
\textbf{ADDITIONAL EXAMPLE 2}
Graph the piecewise function.
(f(x)=\begin{cases}x^2-6,&-4\leq x<0\\8,&0\leq x<5\\-\frac{1}{3}x,&5\leq x<10\end{cases})
% FIG: 40 714 1027 833 1121
\placeholder{graph}{Additional Example 2 orange graph: parabola x^2-6 from filled (-4,10) to open (0,-6); horizontal y=8 from filled (0,8) to open (5,8); line y=-x/3 from filled (5,-5/3) to open (10,-10/3). Labels f(x)=x^2-6, f(x)=8, f(x)=-x/3. x/y arrowed axes, origin O, labels -8,-4,4,8 on x and -8,-4,4,8 on y, grid spacing 2, visible [-10,10] by [-10,10].}
\answer{See graph.}
\textbf{Example 2} Help students determine when function pieces are connected.
Q: How can you determine if two pieces of a function's graph will be connected?
\answer{At the x-coordinate where two intervals of the function pieces come together, if the two functions have the same y-value, then the two pieces will connect.}
\te{Additional Examples are available at SavvasRealize.com. Screenshots show Go back, ESSENTIAL QUESTION, SOLUTION; Example 2 also shows TRY ANOTHER ONE.}
\end{te-addendum}

% Source page 41: 24 STEP 2 Understand & Apply
\begin{tryit}{1}
How much will Alani earn if she works:
a. 37 hours?
b. 43 hours?
\answer{a. \$444 b. \$534}
\end{tryit}
\begin{te-addendum}{2}{GeneralizeETP}
\textbf{EXAMPLE 2 Graph a Piecewise-Defined Function}
\textbf{Use and Connect Mathematical Representations ETP}
Q: Why can x-values only be included in one interval?
\answer{In order to be a function, x-values cannot repeat.}
Q: How do you know the function is decreasing when (-2<x<0)?
\answer{Looking at the graph of the parabola, you can see that it is going down to a minimum when (-2<x<0).}
\end{te-addendum}
\begin{example}{2}
\textbf{Graph a Piecewise-Defined Function}
How do you graph a piecewise defined function?
(f(x)=\begin{cases}4x+11,&-10\leq x<-2\\x^2-1,&-2\leq x\leq2\\2x-10,&2<x<10\end{cases})
What are the domain and range? Over what intervals is the function increasing or decreasing?
% FIG: 41 1008 389 1117 531
\placeholder{graph}{Example 2: red line 4x+11 from filled (-10,-29) to open (-2,3), covered there by blue filled point. Blue parabola x^2-1 from filled (-2,3) through vertex (0,-1) to filled (2,3). Green line 2x-10 from open (2,-6) to open (10,10). x/y axes arrowed, origin O; horizontal grid 2 with labels -8,-4,4,8; vertical grid 3 with labels -24,-18,-12,-6,6. Window approximately x=-10 to 10,y=-30 to 12; no curve arrows.}
\te{Callout: Sketch (y=2x-10) with open circles at (x=2) and (x=10).}
\te{Callout: Sketch (y=x^2-1) with closed circles at (x=\pm2).}
\te{Callout: Sketch (y=4x+11) with a closed circle at (x=-10) and an open circle at (x=-2).}
\te{STUDY TIP: A closed circle on the graph means the coordinates of the point are included in the domain and range of the function. An open circle indicates they are not included.}
To determine the range, calculate the y-values that correspond to the minimum and maximum x-values on the graph. For this graph, these values occur at the endpoints of the domain of the piecewise function, (-10\leq x\leq10).
Evaluate (y=4x+11) for (x=-10).
(y=4(-10)+11)
(y=-29)
Evaluate (y=2x-10) for (x=10).
(y=20-10)
(y=10)
The range is (-29\leq y<10).
The domain is (\{x\mid-10\leq x<10\}). The range is (\{y\mid-29\leq y<10\}). The function is increasing on ((-10,-2)), ((0,2)), and ((2,10)). The function is decreasing on ((-2,0)).
\te{Printed endpoint discussion uses x\leq10, while the rule and final domain use x<10; the latter excludes the endpoint.}
\answer{Domain: (\{x\mid-10\leq x<10\}); range: (\{y\mid-29\leq y<10\}); increasing on ((-10,-2)), ((0,2)), and ((2,10)); decreasing on ((-2,0)).}
\end{example}
\begin{tryit}{2}
Graph the piecewise-defined function. What are the domain and range? Over what intervals is the function increasing or decreasing?
a. (f(x)=\begin{cases}2x+5,&-6\leq x\leq-2\\2x^2-7,&-2<x<1\\-4-x,&1\leq x\leq3\end{cases})
b. (f(x)=\begin{cases}3,&-4<x\leq-1\\-x,&0\leq x\leq2\\3-x,&2<x<4\end{cases})
\answer{a. The domain is (-6\leq x\leq3). The range is (-7\leq y\leq1). The function is increasing when (-6<x<-2) and (0<x<1). The function is decreasing when (-2<x<0) and (1<x<3).}
% FIG: 41 145 464 345 666
\placeholder{graph}{Try It 2a answer: red line from (-6,-7) to (-2,1); blue parabola 2x^2-7 from (-2,1) through vertex (0,-7) to (1,-5); green line from (1,-5) to (3,-7). Continuous joins, no explicit endpoint circles or curve arrows. Unit grid approximately x=-6 to 4,y=-8 to 2; arrowed x/y axes, origin O; x labels -4,-2,2 and y labels -2,-4,-6.}
\answer{b. The domain is (-4<x\leq-1) and (0\leq x<4). The range is (-2\leq y<1) and 3. There are no intervals where the function is increasing. The function is decreasing when (0<x<2) and (2<x<4).}
% FIG: 41 145 757 344 956
\placeholder{graph}{Try It 2b: red horizontal segment from open (-4,3) to filled (-1,3); blue segment from filled (0,0) to filled (2,-2); green segment from open (2,1) to open (4,-1). Unit square grid [-5,5] by [-5,5]; arrowed x/y axes, origin O, labels -4,-2,2,4 on each axis; no curve arrows.}
\end{tryit}
\begin{te-addendum}{2}{ElicitEvidence}
\textbf{Elicit and Use Evidence of Student Thinking ETP}
Q: Are there any gaps in the piecewise-defined functions?
\answer{Yes; in Try It 2b, there is a gap in the domain from (-1) to 0.}
\end{te-addendum}
\begin{te-addendum}{2}{CommonError}
\textbf{Common Error—Try It! 2b}
Some students may state that the domain is (-4\leq x\leq4) and the range is (-2\leq y\leq3). They may not realize the gaps should not be included in the domain or range. Have students run their finger over the graph to check that every number in the domain and range is included.
\end{te-addendum}
\begin{te-addendum}{2}{HabitsOfMind}
\textbf{Use with EXAMPLES 1 \& 2}
Q: Reason Why is the interval for the domain of the second piece of the function in Try It! 2a defined using the (<) symbol rather than the (\leq) symbol?
\answer{The (<) symbol is used since each input can only be matched with one rule and the endpoints ((-2) and 1) are included in the first and third rules.}
\end{te-addendum}


% Source page 42: 25 STEP 2 Understand & Apply
\begin{te-addendum}{3}{GeneralizeETP}
\textbf{EXAMPLE 3 Write a Piecewise-Defined Rule From a Graph}
\textbf{Use and Connect Mathematical Representations ETP}
Q: What do you notice about the graph?
\answer{The graph is composed of three separate linear pieces.}
Q: What does a horizontal line in a graph mean?
\answer{The slope is 0, and there are multiple x-values with the same y-value.}
Q: Given the slope, how do you determine what points to use to find the value of (b) in the slope-intercept form of a linear function (f(x)=mx+b)?
\answer{The x-values can be any values that are located in the domain, and the y-values are the corresponding values in the range.}
\end{te-addendum}
\begin{example}{3}
\textbf{Write a Piecewise-Defined Rule From a Graph}
What is the rule that describes the piecewise-defined function shown in the graph?
% FIG: 42 1035 241 1151 339
\placeholder{graph}{Example 3: blue Piece A line from filled (-2,-10) to open (1,-1); red Piece B line from filled (1,4) to filled (5,8); green Piece C horizontal from (5,8) to filled (8,8), with left endpoint coinciding with filled B endpoint. Color labels A,B,C and point labels (-2,-10),(1,-1),(1,4),(5,8). Arrowed x/y axes, origin O, x labels -8,-4,4,8, y labels -8,-4,4,8; grid spacing 2, window [-10,10] by [-10,10].}
\te{Callout: The values at 1 and 5 are only included in one piece of the graph, so this represents a function.}
There are three pieces for this function. For each piece, find the x interval: ([-2,1)), ([1,5]), and ((5,8]). Since the pieces are linear, find each slope and pick a point to find the rule.
\begin{tabular}{lll}
Piece A & Piece B & Piece C \\
Interval: (-2\leq x<1) & Interval: (1\leq x\leq5) & Interval: (5<x\leq8) \\
point ((-2,-10)), slope (=3) & point ((1,4)), slope (=1) & point ((5,8)), slope (=0) \\
(y=mx+b) & (y=mx+b) & (y=mx+b) \\
(-10=(3)(-2)+b) & (4=(1)(1)+b) & (8=(0)(5)+b) \\
(b=-4) & (b=3) & (b=8) \\
(f(x)=3x-4) & (f(x)=x+3) & (f(x)=8) \\
\end{tabular}
The rule for this function is:
(f(x)=\begin{cases}3x-4,&-2\leq x<1\\x+3,&1\leq x\leq5\\8,&5<x\leq8\end{cases})
\answer{(f(x)=\begin{cases}3x-4,&-2\leq x<1\\x+3,&1\leq x\leq5\\8,&5<x\leq8\end{cases})}
\end{example}
\begin{tryit}{3}
What rule defines the function in each of the following graphs?
a.
% FIG: 42 899 591 1003 692
\placeholder{graph}{Try It 3a: red horizontal segment from filled (-5,7) to filled (-2,7), downward segment from (-2,7) to filled (3,2), separate upward segment from open (4,5) to filled (6,9). Arrowed x/y axes, origin O, labels -8,-4,4,8 on both axes, grid spacing 2; window [-10,10] by [-10,10], no curve arrows.}
\answer{a. (f(x)=\begin{cases}7,&\text{if }-5\leq x\leq-2\\5-x,&\text{if }-2<x\leq3\\2x-3,&\text{if }4<x\leq6\end{cases})}
b.
% FIG: 42 1051 591 1153 692
\placeholder{graph}{Try It 3b: red segment from filled (-3,3) to open (-2,2); horizontal segment from filled (-2,1) to open (1,1); upward segment from filled (1,-1) to filled (2,3). Unit square grid [-5,5] by [-5,5]; arrowed x/y axes, origin O, labels -4,-2,2,4 on both axes; no curve arrows.}
\answer{b. (f(x)=\begin{cases}-x,&\text{if }-3\leq x<-2\\1,&\text{if }-2\leq x<1\\4x-5,&\text{if }1\leq x\leq2\end{cases})}
\end{tryit}
\begin{te-addendum}{3}{ELLAddendum}
\textbf{English Language Learners—Use with EXAMPLE 3}
\textbf{LISTENING—BEGINNING} Ask students to listen as you pronounce the word separate as a verb (sep-uh-rate) and then as an adjective (sep-er-it). Then have students listen to the following phrases and identify whether separate is used as a verb or an adjective: "two separate questions"; "to separate the fighting boys"; "to separate the field by a fence"; "three separate classrooms."
Q: Explain how the word separate is used in this example.
\answer{as an adjective in the phrase separate linear pieces}
Q: When in math might the word separate be used differently than in the example?
\answer{Separate the variable terms from the constants.}
\textbf{SPEAKING—DEVELOPING} In small groups, have students discuss different situations in which they encounter rules.
Q: In the situations you discussed, what happens when a rule is broken?
\answer{A person gets in trouble.}
Q: How are rules used in this example?
\answer{A rule is a way to describe the output for each input in a function.}
Q: Can a mathematical rule be broken?
\answer{No, someone may apply a rule incorrectly, but that would result in a wrong answer, not the rule being broken.}
\textbf{READING—EXPANDING} Distribute a set of cards to groups of students with the terms otherwise, clockwise, and lengthwise on them. Distribute another set of cards with the definitions of those terms. Have students read the words and the definitions and match them together.
Q: What does the suffix -wise seem to indicate?
\answer{a way of doing or being, or in a particular direction}
Q: What does the word piece mean?
\answer{part of a whole}
Q: How can those ideas be combined to give you a better understanding of the term piecewise-defined function?
\answer{a function that is broken up into pieces}
\end{te-addendum}

% Source page 43: 26 STEP 2 Understand & Apply
\begin{te-addendum}{4}{GeneralizeETP}
\textbf{EXAMPLE 4 Rewrite an Absolute Value Function As a Piecewise-Defined Function}
\textbf{Use and Connect Mathematical Representations ETP}
Q: What does ((h,k)) represent in the graph of (f)?
\answer{The point ((h,k)) is the vertex of the graph, where the function has its minimum and where the slope changes sign as (x) increases.}
Q: Why do you choose x-values to determine the slope and equation?
\answer{Choosing x-values on each side of the vertex helps you to determine the slope of each piece. You can then use the points to determine the equation of each line.}
\end{te-addendum}
\begin{example}{4}
\textbf{Rewrite an Absolute Value Function As a Piecewise-Defined Function}
How can you rewrite the function (f(x)=|6x+18|) as a piecewise-defined function?
Step 1 Write the function in the form (f(x)=a|x-h|+k) to find the vertex of the function.
(f(x)=|6x+18|)
(=|6(x+3)|)
(=6|x-(-3)|+0)
\te{Callouts: (h=-3); (k=0).}
The vertex is ((h,k)=(-3,0)). The graph has two linear pieces, one to the left of (x=-3), and one to the right of (x=-3).
Step 2 Graph the function. Determine the slope and equation of each piece by choosing x-values on either side of (-3).
% FIG: 43 911 474 1011 574
\placeholder{graph}{Red V graph f(x)=|6x+18|, filled vertex (-3,0), slopes -6 and 6 with upward arrowheads. Unit grid, arrowed x/y axes, origin O; x labels -8,-6,-4,-2, y labels 2,4,6,8. Window approximately x=-9 to 1,y=-1 to 9; no other marked points.}
\te{Callout: Choose point ((-4,6)). The slope to the vertex is (-6). The equation is (y=-6x-18).}
\te{Callout: Choose point ((-2,6)). The slope to the vertex is 6. The equation is (y=6x+18).}
\te{Callout: The vertex is ((-3,0)).}
\te{GENERALIZE: The parent absolute value function (f(x)=|x|) is a piecewise-defined function: (f(x)=\begin{cases}x,&x\geq0\\-x,&x<0\end{cases}).}
Step 3 Write the piecewise-defined function.
The absolute value function (f(x)=|6x+18|) can be written as the piecewise-defined function:
(f(x)=\begin{cases}-6x-18,&x<-3\\6x+18,&x\geq-3\end{cases})
\answer{(f(x)=\begin{cases}-6x-18,&x<-3\\6x+18,&x\geq-3\end{cases})}
\end{example}
\begin{tryit}{4}
How can you rewrite each function as a piecewise-defined function?
a. (f(x)=|-5x-10|)
b. (f(x)=-|\frac{1}{2}x|+3)
\answer{a. (f(x)=\begin{cases}-5x-10,&\text{if }x<-2\\5x+10,&\text{if }x\geq-2\end{cases}) b. (f(x)=\begin{cases}\frac{x}{2}+3,&\text{if }x<0\\-\frac{x}{2}+3,&\text{if }x\geq0\end{cases})}
\end{tryit}
\begin{te-addendum}{4}{ElicitEvidence}
\textbf{Elicit and Use Evidence of Student Thinking ETP}
Q: How do the piecewise equations compare to the original equation?
\answer{For one piece of the function the absolute value signs are removed, and for the other the addition or subtraction sign is changed.}
\end{te-addendum}
\begin{te-addendum}{4}{HabitsOfMind}
\textbf{Use with EXAMPLES 3, \& 4}
Q: Use Structure Why can the graph of an absolute value function also be defined as a piecewise-defined function?
\answer{Absolute value graphs consist of pieces of two lines with opposite slopes.}
\end{te-addendum}
\begin{te-addendum}{4}{RtISupport}
\textbf{Support Struggling Students}
USE WITH EXAMPLE 4 Some students may struggle with writing a rule for an absolute value function. Have students practice using a graph to write a rule for an absolute value function.
% FIG: 43 777 1066 1032 1320
\placeholder{graph}{Blue V with vertex (8,0), left branch through (0,8) and right branch drawn toward (12,3.2), arrowheads both upper ends; source answer identifies f(x)=|x-8| though drawn right branch is shallower than slope 1. Unit square grid, arrowed x/y axes, origin O, labels 2,4,6,8,10 on each positive axis; window approximately [-1,12] by [-1,12]. No marked points.}
Q: How do you determine the absolute value function from the graph provided? State the function.
\answer{The vertex of the absolute value function is ((8,0)) and the slope is 1, so the function is (f(x)=|x-8|).}
Q: How do you determine a piecewise-defined function from the graph provided? State the function.
\answer{Since this is an absolute value function, the vertex is where the slope changes from negative to positive. The two pieces of the function can be represented as (f(x)=\{-x+8,\text{ if }x<8;\ x-8,\text{ if }x\geq8\}).}
\te{Printed right branch appears shallower than the slope 1 stated in the answer; likely a drawing inconsistency.}
\end{te-addendum}

% Source page 44: 27 STEP 2 Understand & Apply
\begin{te-addendum}{5}{GeneralizeETP}
\textbf{EXAMPLE 5 Graph a Step Function}
\textbf{Use and Connect Mathematical Representations ETP}
Q: What do you notice about the graph?
\answer{An open circle on the graph means the point is not included in that piece of the graph, and a closed circle means the point is included. The graph looks like a staircase.}
Q: How does a step function compare to a piecewise-defined function?
\answer{A step function is a piecewise-defined function containing all constant pieces.}
Q: Are step functions continuous functions? Explain.
\answer{No, step functions are not continuous. You need to remove your pencil from the paper when sketching the graph of a step function, showing that there are gaps.}
\end{te-addendum}
\begin{example}{5}
\textbf{Graph a Step Function}
\textbf{APPLICATION}
The shipping cost of items purchased from an online store is dependent on the weight of the items. The table represents shipping costs (y) based on the weight (x). Graph the function. What are the domain and range of the function, and what do they represent? What are the maximum and minimum values?
\begin{tabular}{lllll}
Weight of Items (lb) & (0<x\leq2) & (2<x\leq4) & (4<x\leq6) & (6<x\leq8) \\
Shipping Cost (\$) & 5 & 8 & 11 & 14 \\
\end{tabular}
The graph of the function looks like the steps of a staircase. This is called a step function since it is a piecewise-defined function consisting of horizontal pieces.
% FIG: 44 817 372 983 514
\placeholder{graph}{Shipping graph over red delivery truck illustration: cyan horizontal steps at y=5 on (0,2], y=8 on (2,4], y=11 on (4,6], y=14 on (6,8]. Each step open left, filled right, no line arrows. x-axis Weight (lb), labels 0,2,4,6,8, grid 2; y-axis Shipping Cost (dollars), labels 2,4,6,8,10,12,14, grid 2. Arrowed x/y axes, origin O, window approximately x=-2 to 10,y=0 to 16.}
\te{Callout: Each y-value goes with an interval of x-values. For example, (y=5) for (0<x\leq2) and (y=8) for (2<x\leq4).}
\te{Callout: A closed circle means that the x-value is included in that domain, and the (\leq) or (\geq) sign is used. An open circle means it is not included.}
Domain: (\{x\mid0<x\leq8\})
Range: (\{5,8,11,14\})
The domain of this function represents all possible weights of the items shipped. The store can ship items that weigh 8 pounds or less. Negative numbers and 0 are not part of the domain because an item must weigh something and cannot have a negative weight.
The range represents all possible costs that a customer could pay to ship the items.
This function has a minimum of 5 and a maximum of 14.
\te{COMMON ERROR: You might think that the range of this function would be the interval [5,14], but only the values 5,8,11, and 14 are possible outputs.}
\answer{Domain: (\{x\mid0<x\leq8\}), possible weights; range: (\{5,8,11,14\}), possible shipping costs; minimum 5, maximum 14.}
\end{example}
\begin{tryit}{5}
The table below represents fees for a parking lot. Graph the function. What are the domain and range of the function, and what do they represent? What are the maximum and minimum values?
\begin{tabular}{lllll}
Time (h) & (0<t\leq3) & (3<t\leq6) & (6<t\leq9) & (9<t\leq12) \\
Cost (\$) & 10 & 15 & 20 & 25 \\
\end{tabular}
% FIG: 44 160 506 442 676
\placeholder{graph}{Try It 5 blue step graph: y=10 on (0,3], y=15 on (3,6], y=20 on (6,9], y=25 on (9,12], all steps open left/filled right. x-axis Time (h), labels 0,2,4,6,8,10,12, grid 1; y-axis Cost (dollars), labels 0,10,20,30, grid 5. Arrowed x/y axes; window [0,12] by [0,30].}
\answer{Domain: (\{x\mid0<x\leq12\}), hours parked in lot. Range: (\{y\mid10,15,20,25\}), cost for parking. Minimum 10; Maximum 25.}
\end{tryit}
\begin{te-addendum}{5}{HabitsOfMind}
\textbf{Use with EXAMPLE 5}
Q: Make Sense and Persevere How would a piecewise-defined rule for the function in the Try It! show that the graph is a step function?
\answer{The rule for each interval would simply be a constant.}
\end{te-addendum}
\begin{te-addendum}{5}{RtIExtend}
\textbf{Extend Student Thinking}
USE WITH EXAMPLE 5 Have students explore creating a table when given the domain and range of the function.
Q: State the minimum and maximum and create a table to represent the following information:
Domain: (\{x\mid-3<x\leq12\})
Range: (\{y\mid-5,4,13,21,30\})
\answer{Answers may vary. Sample: Minimum (= -5); Maximum (=30).}
\begin{tabular}{llllll}
(x) & (-3<x\leq0) & (0<x\leq3) & (3<x\leq6) & (6<x\leq9) & (9<x\leq12) \\
(y) & (-5) & 4 & 13 & 21 & 30 \\
\end{tabular}
\end{te-addendum}


% Source page 45: 28 STEP 2 Understand & Apply
\begin{te-addendum}{ConceptSummary}{GeneralizeETP}
\textbf{CONCEPT SUMMARY Piecewise-Defined Functions}
Q: How can you describe the graph of a piecewise-defined function?
\answer{The graph of a piecewise-defined function has different pieces for different parts in the domain that may not be continuous.}
\te{Concept Summary, Do You Understand, and Do You Know How are available at SavvasRealize.com.}
\end{te-addendum}
\begin{te-addendum}{ConceptSummary}{CommonError}
\textbf{Do You UNDERSTAND? | Do You KNOW HOW?}
\textbf{Common Error—Exercise 9} Students may mistakenly use the wrong symbol to represent the domain, such as (<) instead of (\leq), since they may not pay attention to the open and closed circles. Have students check points on the line to see if they are in the domain they have written.
\end{te-addendum}
\begin{concept-summary}
\textbf{CONCEPT SUMMARY Piecewise-Defined Functions}
\begin{tabular}{ll}
WORDS & ALGEBRA \\
A piecewise function has different rules for intervals of its domain. & (f(x)=\begin{cases}7,&-5\leq x\leq-2\\5-x,&-2<x\leq3\\2x-3,&4<x\leq6\end{cases}) \\
\end{tabular}
% FIG: 45 989 264 1101 366
\placeholder{graph}{Summary graph: red horizontal piece from filled (-5,7) to filled (-2,7); blue descending piece to filled (3,2), open at (-2,7) but covered by red endpoint; green piece from open (4,5) to filled (6,9). All five listed points labeled with color-matched coordinate text. Arrowed x/y axes, origin O, grid 2, labels -8,-4,4,8 on both axes; window [-10,10] by [-10,10]; no curve arrows.}
\textbf{Do You UNDERSTAND?}
1. ESSENTIAL QUESTION How do you model a situation in which different functions describe specific intervals of its domain?
\answer{Model each part of its domain separately, and then combine them into one function using a brace, and include intervals detailing which model goes with each part of the domain.}
2. Vocabulary How do piecewise-defined functions differ from step functions?
\answer{Piecewise functions are two or more functions over different intervals, while step functions are two or more constant values for different intervals.}
3. Error Analysis Given the function
(f(x)=\begin{cases}2x+5,&-2<x\leq4\\-4x-7,&4<x\leq9\end{cases})
Tomás says there is an open circle at (x=4) for both pieces of the function. Explain his error.
\answer{There is an open circle on the piece of the function for (y=-4x-7), but there is a closed circle on the piece of the function for (y=2x+5).}
4. Communicate Precisely What steps do you follow when graphing a piecewise-defined function?
\answer{Divide the graph into the sections defined by the domain. Graph the function given in each section. Note whether the symbols of the domain require an open or closed circle at the edges of each function graph.}
5. Make Sense and Persevere Is the relation defined by the following piecewise rule a function? Explain.
(y=\begin{cases}7x-4,&x<2\\-x+5,&x\geq-2\end{cases})
\answer{No; Sample: Over the interval ([-2,2)) for (x), each input to the rule gives two outputs, so this is not a function.}
\te{Printed "each input" overstates the claim: the two formulas agree at x=9/8; the relation still fails to be a function elsewhere on the overlap.}
\textbf{Do You KNOW HOW?}
Graph the function.
6. (f(x)=\begin{cases}-x+1,&-10\leq x<-3\\x^2-9,&-3\leq x\leq3\\2x+1,&3<x<5\end{cases})
\answer{See graph.}
% FIG: 45 124 941 329 1159
\placeholder{graph}{Answer 6: blue segment from filled (-10,11) to open (-3,4); green parabola x^2-9 from filled (-3,0) through vertex (0,-9) to filled (3,0); red segment from open (3,7) to open (5,11). Grid spacing 2, arrowed x/y axes, origin O; x labels -8,-4,4,8 and y labels -8,-4,4,8. Window [-10,10] by [-10,12]; no curve arrows.}
7. (g(x)=\begin{cases}1,&0\leq x<2\\3,&2\leq x<4\\5,&4\leq x<6\\7,&6\leq x<8\end{cases})
\answer{See graph.}
% FIG: 45 124 1178 320 1372
\placeholder{graph}{Answer 7: blue horizontal steps y=1 on [0,2), y=3 on [2,4), y=5 on [4,6), y=7 on [6,8). Filled left/open right endpoints, no curve arrows. Unit grid [0,9] by [0,9]; arrowed x/y axes, labels 0,2,4,6,8 on both axes.}
8. Given the function
(f(x)=\begin{cases}-2x+4,&0\leq x<8\\-5x+11,&x\geq8\end{cases})
is the function increasing or decreasing over the interval ([2,7])? Find the rate of change over this interval.
\answer{decreasing; (-2)}
9. What is the rule that defines the function shown in the graph?
% FIG: 45 916 631 1021 732
\placeholder{graph}{Question 9: red horizontal steps y=3 from open (-2,3) to filled (0,3); y=1/2 from open (0,1/2) to filled (2,1/2); y=-2 from open (2,-2) to filled (4,-2). Unit square grid [-5,5] by [-5,5], arrowed x/y axes, origin O; labels -4,-2,2,4 on each axis, no curve arrows.}
\answer{(f(x)=\begin{cases}3,&\text{if }-2<x\leq0\\\frac{1}{2},&\text{if }0<x\leq2\\-2,&\text{if }2<x\leq4\end{cases})}
\end{concept-summary}


% Source page 46: 29 STEP 3 Practice & Problem Solving
\begin{te-addendum}{Practice}{GeneralizeETP}
\textbf{PRACTICE \& PROBLEM SOLVING}
\textbf{Lesson Practice} You may opt to have students complete the automatically scored Practice and Problem Solving items online powered by MathXL for School.
Choose from: Lesson Practice; Adaptive Practice.
You may also take advantage of the bank of exercises for assigning additional practice.
\textbf{Assignment Guide}
\begin{tabular}{ll}
On-Level & Advanced \\
10--12, 15--27 & 10--27 \\
\end{tabular}
\textbf{Engage Through Student Choice}
Promote student agency by allowing students to choose practice items. You may structure this choice in many ways.
For example:
Assign each section a point value. Students choose at least one item from each section and items chosen should have a minimum of 20 total points.
\begin{tabular}{ll}
Understand, Apply & 2 points each \\
Practice & 1 point each \\
Assessment Practice & 1 point each \\
Performance Task & 3 points \\
\end{tabular}
\te{Practice \& Problem Solving and video tutorials are available at SavvasRealize.com. Scan the QR code to access a video tutorial. Printed answers for Exercises 16--21 appear on PDF page 47 and are attached below.}
\end{te-addendum}
\begin{item-analysis}
\begin{tabular}{lll}
\toprule
Example & Items & DOK \\
\midrule
1 & 15, 22, 23 & 2 \\
2 & 12, 13, 16, 25 & 1 \\
2 & 11 & 2 \\
3 & 17 & 1 \\
3 & 27 & 3 \\
4 & 18, 19 & 1 \\
4 & 26 & 2 \\
5 & 14, 20, 21 & 1 \\
5 & 10, 24 & 2 \\
\bottomrule
\end{tabular}
\end{item-analysis}
\begin{practice}{10}{2}
UNDERSTAND. Communicate Precisely What do closed circles and open circles on the graph of a step function indicate?
\answer{Closed circles indicate that the values of a point are included in the domain and range of a function. Open circles indicate that the values are not included.}
\end{practice}
\begin{practice}{11}{2}
Error Analysis What error did Mateo make when defining the domain of the graph? Explain.
% FIG: 46 548 386 757 593
\placeholder{graph}{Error Analysis card: red shallow increasing ray extending left from open (3,-1), through approximately (0,-2) and (-3,-3); blue ray from open (3,2) through (5,4) extending upper-right. Unit grid, arrowed x/y axes, origin O, labels -4,-2,2,4 on x and -4,-2,2,4 on y. Window approximately [-6,6] by [-6,6]. Text "domain: all real numbers" and large red X at lower right.}
\answer{The graph has open circles at (x=3), so the domain is all real numbers except 3.}
\end{practice}
\begin{practice}{12}{1}
Communicate Precisely For what values of (x) is the function (f(x)=\begin{cases}-3x+4,&-2<x\leq3\\2x+1,&4\leq x<9\end{cases}) defined?
\answer{(-2<x\leq3) and (4\leq x<9)}
\end{practice}
\begin{practice}{13}{1}
Mathematical Connections For the piecewise-defined function (f(x)=\begin{cases}7,&x>3\\5x-3,&x\leq3\end{cases}), find two x-values that have the same y-value and the sum of the x-values is 10.
\answer{Answers may vary. Samples: (x=2) and (x=8).}
\end{practice}
\begin{practice}{14}{1}
Higher Order Thinking The function (f(x)=\lfloor x\rfloor) is called the greatest integer function because the output returned is the greatest integer less than or equal to (x). For example, (f(3.2)=\lfloor3.2\rfloor=3) and (f(0.975)=\lfloor0.975\rfloor=0). Graph the function (f(x)=\lfloor x\rfloor). What type of graph does this look like?
\answer{step function}
\te{No answer graph is printed for this item.}
\end{practice}
\begin{practice}{15}{2}
PRACTICE. A cell phone company restricts the data speeds for customers after they have used 10 GB of data in a month. Download speeds drop from a guaranteed rate of 43 Mbps to 36 Mbps. Write a piecewise-defined function (T(x)), to compute the time in minutes to download 1 GB of data no matter if it is the 5th GB or the 15th GB. Include in your function the proper conversion using (8000\text{ Mb}\approx1\text{ GB}). SEE EXAMPLE 1
\answer{(T(x)=\begin{cases}\frac{x}{43(60)},&x\leq80{,}000\\\frac{x-80000}{36(60)}+31.01,&x>80{,}000\end{cases})}
\te{Printed answer models cumulative download time for x megabits; the prompt asks about the time to download a single GB. The printed function is retained.}
\end{practice}
\begin{practice}{16}{1}
Graph the piecewise-defined function. State the domain and range. Identify whether the function is increasing, constant, or decreasing on each interval of the domain. SEE EXAMPLE 2
(f(x)=\begin{cases}\frac{1}{4}x+3,&-2<x\leq0\\2,&0<x\leq4\\3-x,&4<x\leq7\end{cases})
\answer{domain: (-2<x\leq7); range: (-4\leq y<-1) and (y=2) and (2.5<y\leq3); increasing: between (-2) and 0; constant: between 0 and 4; decreasing: between 4 and 7.}
% FIG: 47 126 164 324 359
\placeholder{graph}{Answer 16: red line from open (-2,2.5) to filled (0,3); blue horizontal segment from open (0,2) to filled (4,2); green line from open (4,-1) to filled (7,-4). Unit grid, arrowed x/y axes, origin O; x labels -2,2,4,6 and y labels -4,-2,2,4. Window x=-3 to 7,y=-5 to 5; no curve arrows.}
\end{practice}
\begin{practice}{17}{1}
Write the rule that defines the function in the following graph. SEE EXAMPLE 3
% FIG: 46 855 589 994 723
\placeholder{graph}{Red ray y=x+5 ending at filled (-1,4) and extending southwest through (-5,0), arrowhead at lower-left; red horizontal ray y=-2 beginning at open (1,-2) and extending right. Unit grid [-5,5] by [-5,5], arrowed x/y axes, origin O, labels -4,-2,2,4 on both axes.}
\answer{(y=\begin{cases}x+5,&\text{when }x\leq-1\\-2,&\text{when }x>1\end{cases})}
\end{practice}
\begin{practice}{18}{1}
Write each absolute value function as a piecewise-defined function. SEE EXAMPLE 4
(f(x)=|3x+1|)
\answer{(f(x)=\begin{cases}3x+1,&\text{if }x\geq-\frac{1}{3}\\-3x-1,&\text{if }x<-\frac{1}{3}\end{cases})}
\end{practice}
\begin{practice}{19}{1}
Write each absolute value function as a piecewise-defined function. SEE EXAMPLE 4
(g(x)=|-2x-6|)
\answer{(f(x)=\begin{cases}-2x-6,&\text{if }x<-3\\2x+6,&\text{if }x\geq-3\end{cases})}
\te{Printed answer uses f(x), while the prompt names g(x).}
\end{practice}
\begin{practice}{20}{1}
Graph the step function. SEE EXAMPLE 5
(f(x)=\begin{cases}2,&-3\leq x<1\\5,&1\leq x<4\\8,&4\leq x<6\\9,&6\leq x<10\end{cases})
\answer{See graph.}
% FIG: 47 124 638 321 835
\placeholder{graph}{Answer 20: blue horizontal steps y=2 on [-3,1), y=5 on [1,4), y=8 on [4,6), y=9 on [6,10). Filled left/open right endpoints, no line arrows. Grid spacing 2, arrowed x/y axes, origin O; x labels -8,-4,4,8 and y labels -8,-4,4,8; window [-10,10] by [-10,10].}
\end{practice}
\begin{practice}{21}{1}
The parking rates for a parking garage are shown. Graph the function for the cost of parking rates at the garage. SEE EXAMPLE 5
% FIG: 46 990 866 1122 1044
\placeholder{illustration}{Parking Rates sign: 4 dollars per half hour; 20 dollars maximum for 12 hours; red heading and dark-blue rate panels with car symbol.}
\answer{See graph.}
% FIG: 47 123 858 424 1120
\placeholder{graph}{Answer 21, Parking Cost: blue horizontal steps y=4 on (0,0.5], y=8 on (0.5,1], y=12 on (1,1.5], y=16 on (1.5,2], y=20 on (2,12]. Open left/filled right endpoints. x-axis Hours, labels 0,2,4,6,8,10,12 with grid 1; y-axis Cost (dollars), labels 0,4,8,12,16,20 with grid 2. Arrowed x/y axes, no line arrows.}
\end{practice}

% Source page 47: 30 STEP 3 Practice & Problem Solving
\begin{te-addendum}{Practice}{GeneralizeETP}
\textbf{Answers}
\te{The answers printed here for Exercises 16--21 are attached to those practice blocks above. Practice \& Problem Solving is available at SavvasRealize.com.}
\end{te-addendum}
\begin{practice}{22}{2}
APPLY. Model With Mathematics If Keandre works more than 40 h per week, his hourly wage for the extra hour(s) is 1.5 times the normal hourly wage of \$12 per hour. Write a piecewise-defined function that gives Keandre's weekly pay (P) in terms of the number (h) of hours he works. Determine how much Keandre will get paid if he works 45 h.
\answer{Sample answer: (P(h)=\begin{cases}12h,&0<h\leq40\\18h-240,&40<h\end{cases}); \$570.}
\end{practice}
\begin{practice}{23}{2}
Model With Mathematics Data plans offered at a phone company, along with overage charges, are shown. Fees are per month.
% FIG: 47 520 463 777 584
\placeholder{illustration}{Four phone-shaped data-plan cards. BASIC Data Plan: 15 dollars for 2 GB, then 15 dollars for each additional GB. GOOD Data Plan: 25 dollars for 4 GB, then 20 dollars for each additional GB. BETTER Data Plan: 40 dollars for 5 GB, then 25 dollars for each additional GB. UNLIMITED Data Plan: 80 dollars.}
a. Write a function for each plan where (x) is the amount of data and (f(x)) is the total monthly cost.
\answer{a. Basic Data Plan (f(x)=\begin{cases}15,&x\leq2\\15+15(x-2),&x>2\end{cases}). Good Data Plan (f(x)=\begin{cases}25,&x\leq4\\25+20(x-4),&x>4\end{cases}). Better Data Plan (f(x)=\begin{cases}40,&x\leq5\\40+25(x-5),&x>5\end{cases}). Unlimited Data Plan (f(x)=80).}
b. Jaheem uses approximately 6 GB of data per month. What is the monthly cost under each data plan?
\answer{b. Basic (f(6)=\$75); Good (f(6)=\$65); Better (f(6)=\$65); Unlimited (f(6)=\$80).}
c. Write an interval for the amount of data that would make each plan the best one to purchase.
\answer{c. Basic ([0,2]); Good ((2,4]); Better ((4,6]); Unlimited ((6,\infty)).}
\te{Printed best-plan intervals do not match all cost crossover points of the printed functions; retained without silently replacing them.}
\end{practice}
\begin{practice}{24}{2}
Reason The cost (C) (in dollars) of sending next-day mail depends on the weight (x) (in ounces) of a package. The cost of packages, up to 5 lb, is given by the function below. What are the domain and range of the function?
(f(x)=\begin{cases}12.25,&0<x\leq8\\16.75,&8<x\leq32\\19.50,&32<x\leq48\\23.50,&48<x\leq64\\25.25,&64<x\leq80\end{cases})
\answer{domain: (\{x\mid0<x\leq80\}); range: (\{y\mid12.25,16.75,19.50,23.50,25.25\}).}
\end{practice}
\begin{practice}{25}{1}
ASSESSMENT PRACTICE. Does the function have a range of ((-\infty,4))? Write yes or no.
a. (f(x)=\begin{cases}x-3,&\text{if }x<-2\\5-x,&\text{if }x>1\end{cases})
b. (h(x)=\begin{cases}x-1,&\text{if }x<-1\\3-x,&\text{if }x>2\end{cases})
c. (g(x)=\begin{cases}x+5,&\text{if }x<-1\\-x-5,&\text{if }x>1\end{cases})
d. (k(x)=\begin{cases}x+2,&\text{if }x<-2\\x+4,&\text{if }x>1\end{cases})
\answer{a. Yes; b. No; c. Yes; d. No.}
\te{Answers are the marked magenta cells in the printed Yes/No table.}
\end{practice}
\begin{practice}{26}{2}
SAT/ACT What is the vertex of the absolute value function (f(x)=-|x-a|+b) where (a) and (b) are real numbers?
A. ((a,b))
B. ((-a,b))
C. ((a,-b))
D. ((-a,-b))
\answer{A}
\end{practice}
\begin{practice}{27}{3}
Performance Task Yama works a varying number of hours per month for a construction company. The following scatter plot shows how much money he earns for each number of hours he works. Write the piecewise-defined function that represents Yama's earnings as a function of his hours worked.
% FIG: 47 824 702 993 922
\placeholder{graph}{Yama's earnings: blue plotted filled points (10,200),(20,400),(30,600),(40,800),(50,1100),(60,1400),(70,1700),(80,2100),(90,2500),(100,2900), joined by blue straight segments beginning at (0,0). x-axis Hours worked, labels 0,20,40,60,80,100, grid 10; y-axis Money Earned (in dollars), labels 0,400,...,3200, grid 200. Arrowed x/y axes; visible [0,110] by [0,3400], no curve arrows.}
\answer{(y=\begin{cases}20x,&0\leq x\leq40\\30x-400,&40\leq x\leq70\\40x-1{,}100,&70\leq x\leq100\end{cases})}
\te{Printed intervals share endpoints 40 and 70; the adjacent formulas agree there.}
\end{practice}


% Source page 48: 30A STEP 4 Assess & Differentiate
\begin{te-addendum}{Assess}{RtISupport}
\textbf{LESSON QUIZ}
Use the Lesson Quiz to assess students' understanding of the mathematics in the lesson.
Students can take the Lesson Quiz online or you can download a printable copy from SavvasRealize.com. The Lesson Quiz is also available in the Assessment Sourcebook.
\textbf{Item Analysis}
\begin{tabular}{lll}
Item & DOK & Skills Review and Practice \\
1 & 2 & F21 \\
2 & 2 & F21 \\
3 & 2 & F21 \\
4 & 2 & F21 \\
5 & 2 & F21 \\
\end{tabular}
Use the student scores on the Lesson Quiz to prescribe differentiated assignments.
If students take the Lesson Quiz online, it will be automatically scored and appropriate differentiated practice will be assigned based on student performance.
\begin{tabular}{lll}
Intervention & 0--3 points & Reteach to Build Understanding; Mathematical Literacy and Vocabulary; Lesson Virtual Nerd videos \\
On-Level & 4 points & Enrichment \\
Advanced & 5 points & Enrichment \\
\end{tabular}
\end{te-addendum}
\begin{te-addendum}{Assess}{GeneralizeETP}
\textbf{ASSESSMENT RESOURCES}
\textbf{1-3 Lesson Quiz: Piecewise-Defined Functions}
Name blank.
1. Which rule defines the function in the graph?
% FIG: 48 837 310 937 412
\placeholder{graph}{Quiz 1 black segments: y=x-1 from (-3,-4) to open (1,0), y=-3x+9 from filled (2,3) to (4,-3). Outer endpoints lack distinct circles. Unit square grid [-5,5] by [-5,5], arrowed x/y axes, origin O, labels -4,-2,2,4 on both axes. No curve arrows.}
A. (f(x)=\begin{cases}x-1,&\text{if }-3<x\leq1\\-3x+9,&\text{if }2<x<4\end{cases})
B. (f(x)=\begin{cases}x-1,&\text{if }-3\leq x<1\\-3x+9,&\text{if }2\leq x\leq4\end{cases})
C. (f(x)=\begin{cases}x-1,&\text{if }-3\leq x<1\\x-3,&\text{if }2\leq x\leq4\end{cases})
D. (f(x)=\begin{cases}x-1,&\text{if }-3<x\leq1\\x-3,&\text{if }2<x<4\end{cases})
\answer{1. B}
2. A scuba diver returning to the surface swims to move up quickly, then waits a few meters below the surface to avoid pressure issues. They complete the ascent more slowly. The function (f) represents the diver's height above sea level over time, in seconds. Which of the following is true for (f)?
(f(x)=\begin{cases}3x+1,&-9<x\leq-2\\-5,&-2<x\leq1\\x-6,&1<x<7\end{cases})
A. The domain is the set of times (-9<x<-2), (-2<x<1), and (1<x<7).
B. The domain is the depth interval ((-26,1)).
C. The domain is the time interval ((-9,7)).
D. The domain is the depth interval ((-9,1)).
\answer{2. C}
3. Write the piecewise-defined function shown in the graph.
% FIG: 48 1050 309 1153 439
\placeholder{graph}{Quiz 3: black V with vertex (9,0), left branch through (0,27), right branch through (18,27), both upper ends arrowed. x/y axes arrowed, origin O, grid spacing 3 on both axes; x labels 6,12,18; y labels 6,12,18,24,30. Window approximately x=-3 to 24,y=-3 to 36, no point circles.}
(f(x)=\begin{cases}\text{blank}(x+\text{blank}),&x\leq\text{blank}\\\text{blank}(x+\text{blank}),&x>\text{blank}\end{cases})
\answer{3. (f(x)=\begin{cases}-3(x+(-9)),&x\leq9\\3(x+(-9)),&x>9\end{cases})}
4. The table shows the admission prices for an amusement park. Graph the function.
\begin{tabular}{ll}
Age & Admission Price \\
under 5 & free \\
5--17 & \$25 \\
18--64 & \$40 \\
65 and over & \$35 \\
\end{tabular}
\answer{4. See graph.}
% FIG: 48 967 631 1121 746
\placeholder{graph}{Quiz 4 magenta step graph: y=0 on (0,5), open circles at (0,0),(5,0); y=25 from filled (5,25) to open (18,25); y=40 from filled (18,40) to open (65,40); y=35 from filled (65,35) with rightward arrow. x-axis Age (years), labels 0,10,20,30,40,50,60, grid 5; y-axis Admission price (dollars), labels 0,10,20,30,40, grid 5. Arrowed x/y axes, origin O; visible x=-5 to 70,y=-5 to 45.}
5. Identify the domain and range of the function in Item 4. Write the range values in ascending order.
domain: (x>\text{blank})
range: (\{\text{blank},\text{blank},\text{blank},\text{blank}\})
\answer{5. domain: (x>0); range: (\{0,25,35,40\}).}
\te{Lesson Quiz is available at SavvasRealize.com. enVision Algebra 2 • Assessment Sourcebook. Copyright © Savvas Learning Company LLC. All Rights Reserved.}
\end{te-addendum}

% Source page 49: 30B STEP 4 Assess & Differentiate
\begin{te-addendum}{Assess}{GeneralizeETP}
\textbf{DIFFERENTIATED RESOURCES}
I = Intervention; O = On-Level; A = Advanced.
\textbf{Reteach to Build Understanding—I}
Provides scaffolded reteaching for the key lesson concepts.
MathXL for School.
\textbf{1-3 Reteach to Build Understanding: Piecewise-Defined Functions}
Name blank.
A piecewise-defined function has different rules for different parts, or pieces, of its domain. In order to evaluate a piecewise-defined function for a given value of (x), find the interval, or piece, that (x) belongs to. Find the corresponding definition of the function for that interval.
1. Given the function
(f(x)=\begin{cases}3,&-5\leq x<-2\\x-2,&-2\leq x\leq0\end{cases})
Solve the piecewise-defined functions for the values listed. Use the equations shown for each group.
Use (f(x)=3), (-5\leq x<-2) to solve.
a. (f(-5)=\text{blank}) \answer{3}
b. (f(-4)=\text{blank}) \answer{3}
c. (f(-3)=\text{blank}) \answer{3}
Use (f(x)=x-2), (-2\leq x\leq0) to solve.
d. (f(-2)=\text{blank}) \answer{(-4)}
e. (f(-1)=\text{blank}) \answer{(-3)}
f. (f(0)=\text{blank}) \answer{(-2)}
2. Deon graphed the piecewise-defined function. Explain what mistake he made and how he should fix it.
(f(x)=\begin{cases}2,&-4\leq x\leq-2\\2x-4,&-1\leq x<2\\3x,&2\leq x<5\end{cases})
% FIG: 49 337 569 403 635
\placeholder{graph}{Reteach 2 incorrect graph: black horizontal segment y=2 from filled (-4,2) to filled (-2,2); black line y=2x-4 from (-1,-6) toward open (5,6). Unit square grid approximately [-5,5] by [-8,10], arrowed x/y axes, origin O; x labels -4,-2,2,4; y labels -8,-4,4,8. Lower left endpoint appears filled; open upper right.}
\answer{Sample answer: Deon did not use the correct domains when graphing the function. He used (2x-4) over the interval (-1<x<5). He should have used (3x) over the domain (2\leq x<5).}
\te{Printed explanation uses a strict lower bound -1<x, whereas the displayed second rule includes x=-1.}
3. Use the piecewise-defined function to fill in the blanks.
(f(x)=\begin{cases}4,&-4<x<-2\\2x-4,&-1<x<2\\3x,&2\leq x<5\end{cases})
a. The domain blank is used when graphing the function (f(x)=2x-4).
\answer{(-1<x<2)}
b. The equation blank is used to find (f(4)=12).
\answer{(f(x)=3x)}
\end{te-addendum}
\begin{te-addendum}{Assess}{GeneralizeETP}
\textbf{Additional Practice—I, O}
Provides extra practice for each lesson.
MathXL for School.
\textbf{1-3 Additional Practice: Piecewise-Defined Functions}
Name blank.
1. A phone company offers a monthly data plan for \$10 a month. The plan includes 2 megabytes of data, and charges \$0.10 per megabytes above the 2 megabytes of data. Write a piecewise-defined function for (M(x)), the cost for (x) megabytes of data used in a month.
\answer{(M(x)=\begin{cases}10,&0<x\leq2\\10+0.10x,&x>2\end{cases})}
\te{Printed second branch charges for all x megabytes; likely intended 10+0.10(x-2) for the stated overage policy.}
2. Graph the piecewise-defined function. State the domain and range. Identify whether the function is increasing, constant, or decreasing on each interval of the domain.
(f(x)=\begin{cases}2,&-4\leq x\leq-2\\x+2,&-2<x<3\\-3x+12,&3\leq x\leq5\end{cases})
\answer{domain: (-4\leq x\leq5); range: (-3\leq y<5); constant when (-4\leq x\leq-2); increasing when (-2<x<3); decreasing when (3\leq x\leq5).}
% FIG: 49 709 465 776 532
\placeholder{graph}{Additional Practice 2 magenta graph: horizontal segment from filled (-4,2) to filled (-2,2); line x+2 from open (-2,0) to open (3,5); line -3x+12 from filled (3,3) to filled (5,-3). Unit grid approximately x=-5 to 5,y=-4 to 6; arrowed x/y axes, origin O; x labels -2,2,4, y labels -2,2,4, no curve arrows.}
3. Write the rule that defines the function in the graph.
% FIG: 49 709 546 776 614
\placeholder{graph}{Additional Practice 3 black graph: horizontal y=6 ray extending left from open (0,6); line y=4x-6 from filled (0,-6) to filled (4,10). Arrowed x/y axes, origin O; unit horizontal grid, vertical grid 2; labels -4,-2,2,4 on x and -4,4,8 on y. Window approximately [-5,5] by [-8,12]; left horizontal arrow, no arrows on rising segment.}
\answer{(f(x)=\begin{cases}6,&x<0\\4x-6,&0\leq x\leq4\end{cases})}
4. Write the function as a piecewise-defined function.
(f(x)=|2x-8|)
\answer{(f(x)=\begin{cases}-2x+8,&x<4\\2x-8,&x\geq4\end{cases})}
5. A shipping service uses the weight of a package to determine its postage. The charge is \$3 for the first pound and \$2 for each additional pound up to 5 pounds. What are the domain and range of the function?
(f(x)=\begin{cases}3,&0<x\leq1\\5,&1<x\leq2\\7,&2<x\leq3\\9,&3<x\leq4\\11,&4<x\leq5\end{cases})
\answer{Domain: (0<x<5). Range: (\{3,5,7,9,11\}).}
\te{Printed answer excludes 5 although the displayed rule includes it; likely 0<x\leq5.}
6. You plan to rent a car from XYZ Car Rental Company for a flat rate of \$35 a day. If you plan to use the car for 3 days or fewer, you must also pay a \$10 insurance fee per day. If you plan to use the car for more than 3 days, there is a \$5 insurance fee per day. Write a piecewise-defined function that models this function.
\answer{(f(x)=\begin{cases}35+10x,&1\leq x\leq3\\35+5x,&x>3\end{cases})}
\te{Printed answer treats the 35-dollar daily rental rate as a one-time amount; likely 35x should replace 35 in both branches.}
\end{te-addendum}
\begin{te-addendum}{Assess}{GeneralizeETP}
\textbf{Enrichment—O, A}
Presents engaging problems and activities that extend the lesson concepts.
MathXL for School.
\textbf{1-3 Enrichment: Piecewise-Defined Functions}
Name blank.
The following system of equations describes the graph of the piecewise-function.
(f(x)=\begin{cases}2x+200,&-75\leq x\leq-60\\x-15,&-60<x\leq-50\\\uncertain{|\frac{3}{2}x|}{coefficient or absolute-value marks unclear in scan},&-50<x<-38\\4x+50,&-38\leq x<-2\\x^2-4,&-2\leq x\leq8\\x^2+6x-5,&8<x\leq15\end{cases})
% FIG: 49 1000 430 1064 492
\placeholder{graph}{Enrichment original black graph on grid approximately x=-80 to 20,y=-120 to 120. Arrowed x/y axes with x labels -60,-40,-20 and y labels -120,60,120; grid spacing approximately 20 horizontally and 30 vertically. Separate upper-left rising segment near (-75,50) to (-60,80), curved segment near (-50,75) to open (-38,57), lower rising segment near (-38,-102) to (-2,42), upward parabola near origin, and steep descending right-hand curve. Tiny endpoints and the correspondence of the right-hand curve to the printed positive quadratic are \uncertain{not clearly resolved}{source is low-resolution}.}
Connect the system of equations of the piecewise-function by rewriting the system so that the points connect. Then graph the new system. You may alter each function by translating it up or down, reflecting it, compressing it, and/or stretching it. Your graph must include the following turning points: ((-75,50)), ((-60,80)), ((-50,-65)), ((-38,-102)), ((-2,0)), ((8,60)), and ((15,-140)).
% FIG: 49 851 555 919 622
\placeholder{graph}{Enrichment answer graph: connected black pieces through (-75,50),(-60,80),(-50,-65),(-38,-102),(-2,0),(8,60),(15,-140); first four and last pieces linear, piece on [-2,8] follows y=x^2-4 with vertex (0,-4). Small endpoint circles not clearly distinguishable. Arrowed x/y axes, origin O; grid approximately 20 horizontally and 30 vertically, x labels -60,-40,-20, y labels 120,60,-120; window approximately x=-80 to 20,y=-150 to 150.}
\answer{(f(x)=\begin{cases}2x+200,&-75\leq x\leq-60\\-\frac{29}{2}x-790,&-60<x\leq-50\\-\frac{37}{12}x-\frac{1315}{6},&-50<x<-38\\\frac{17}{6}x+\frac{17}{3},&-38\leq x<-2\\x^2-4,&-2\leq x\leq8\\-\frac{200}{7}x+\frac{2020}{7},&8<x\leq15\end{cases})}
\te{Printed original last rule has a positive leading quadratic term, while the small graph descends at the right; the scan's tiny original third expression is marked uncertain.}
\end{te-addendum}
\begin{te-addendum}{Assess}{ELLAddendum}
\textbf{Mathematical Literacy and Vocabulary—I, O}
Helps students develop and reinforce understanding of key terms and concepts.
\textbf{1-3 Mathematical Literacy and Vocabulary: Piecewise-Defined Functions}
Name blank.
Name the steps that are used to graph the piecewise-defined function.
(f(x)=\begin{cases}3,&-5\leq x<-3\\x+4,&-3\leq x\leq4\end{cases})
\begin{tabular}{ll}
(-5\leq x<-3) & Choose the first domain. \\
(y=3) & Graph the function on the first domain. \\
(-3\leq x\leq4) & Choose the second domain. \\
(y=x+4) & Graph the function on the second domain. \\
\end{tabular}
% FIG: 49 175 1086 246 1152
\placeholder{graph}{Vocabulary model black graph: y=3 from filled (-5,3) to open (-3,3); y=x+4 from filled (-3,1) to filled (4,8). Unit grid approximately x=-6 to 5,y=-1 to 9; arrowed x/y axes, origin O; x labels -4,-2,2,4 and y labels 2,4,6,8. No line arrows.}
Write the steps that are used to graph the piecewise-defined function.
(f(x)=\begin{cases}2,&-4\leq x\leq-2\\x+4,&-1\leq x\leq4\end{cases})
blank. Choose the first domain.
\answer{(-4\leq x\leq-2)}
blank. Graph the function on the first domain.
\answer{(y=2)}
blank. Choose the second domain.
\answer{(-1\leq x\leq4)}
blank. Graph the function on the second domain.
\answer{(y=x+4)}
% FIG: 49 175 1240 242 1309
\placeholder{graph}{Vocabulary answer magenta graph: horizontal y=2 from filled (-4,2) to filled (-2,2), and y=x+4 from filled (-1,3) to filled (4,8). Unit grid approximately x=-5 to 5,y=-1 to 9; arrowed x/y axes, origin O; x labels -4,-2,2,4 and y labels 2,4,6,8. No line arrows.}
\end{te-addendum}
\begin{te-addendum}{Assess}{GeneralizeETP}
\textbf{Digital Differentiated Resources}
The Reteach to Build Understanding, Additional Practice, and Enrichment activities are available as digital assignments. These activities are automatically assigned when students complete the lesson quiz online and are automatically scored.
\textbf{Digital activity screenshot}
Solve the system of equations by graphing.
(y=3x+4)
(y=-2x+4)
Use the graphing tool to graph the system.
% FIG: 49 813 1047 934 1158
\placeholder{graph}{Digital exercise blank Cartesian grid, arrowed x/y axes, origin O, unit grid, even-number labels -10 through 10 on both axes. Window [-10,10] by [-10,10]; no plotted points or lines; Question Help menu overlays upper-right quadrant.}
Click to enlarge graph.
Click the graph, choose a tool in the palette and follow the instructions to create your graph.
\te{Interface labels: Exit; Question Help; Help Me Solve This; View an Example; Video; Textbook; Glossary; Math Tools; Print; 1 part remaining; Clear All; Check Answer; Review progress; Question 5 of 14; Go; Back; Next. Resource previews carry SavvasRealize.com, enVision Algebra 2 • Teaching Resources, and a copyright notice.}
\end{te-addendum}

\end{document}
